% GENERATED FILE. Edit canonical lessons, then run npm run generate:books.
% canonical-source-hash: fnv1a64:7b15b3a959b01aba
% canonical-lessons: FR-C451-retrait, FR-C451-crediter, FR-C451-regler, FR-C451-decouvert, FR-C451-releve

\chapter{À la banque}
\label{ch:fr-a-la-banque}
\hlchaptermodality{451}

\begin{chapteropening}
\textbf{By the end of this chapter:} \emph{I can deal with a bank account: the withdrawal, paying money in, settling a bill, the overdraft and the statement.}

Complete the last lesson of chapter 451: I can deal with a bank account: the withdrawal, paying money in, settling a bill, the overdraft and the statement.
\end{chapteropening}

\section[le retrait]{le retrait — the withdrawal}

\label{lesson:FR-C451-retrait}



\begin{quote}
Before the new one: say the French for the card machine, then the French for cash.
\end{quote}

\subsection*{You'll want to know: le retrait}
\textbf{le retrait} — \textquotedblleft{}the withdrawal\textquotedblright{}.

\begin{quote}
Il a fait un retrait de cinquante euros.
\end{quote}

\begin{quote}
\emph{He took out fifty euros.}
\end{quote}

\emph{Un retrait d'argent} — a cash withdrawal.

\begin{cousinweb}
\emph{retrait} is the old past participle of \emph{retraire} 'to draw back', from Latin \emph{retrahere}, built on \emph{trahere} 'to draw' (Lewis \& Short). EtymDB does not list it.
\end{cousinweb}

\begin{grammarlens}[title={What you've built}]
One more word for telling, explaining and describing at B1.
\end{grammarlens}

\subsection*{Your turn}
\begin{itemize}
\raggedright
  \item \emph{Say it:} \emph{le retrait}
  \item \emph{Say it:} \emph{le retrait}, once more
  \item \emph{Say it:} a sentence of your own with \emph{le retrait}
  \item \emph{Recall:} say the French for the card machine, then the French for cash, then say \emph{le retrait} again
\end{itemize}

\begin{tcolorbox}[breakable,colback=teal!4,colframe=teal!35!black,title={Before you move on}]
What does le retrait mean? (\textquotedblleft{}The withdrawal\textquotedblright{}.) Say it once more.
\end{tcolorbox}
\section[créditer]{créditer — to credit, to pay into}

\label{lesson:FR-C451-crediter}



\begin{quote}
Before the new one: say the French for cash, then the French for the withdrawal.
\end{quote}

\subsection*{You'll want to know: créditer}
\textbf{créditer} — \textquotedblleft{}to credit, to pay into\textquotedblright{}.

\begin{quote}
La banque a crédité mon compte ce matin.
\end{quote}

\begin{quote}
\emph{The bank paid the money into my account this morning.}
\end{quote}

The opposite of \emph{débiter}.

\begin{cousinweb}
\emph{créditer} is built on \emph{crédit}, from Latin \emph{crēditum} 'a loan', from \emph{crēdere} 'to trust' (Lewis \& Short).
\end{cousinweb}

\begin{grammarlens}[title={What you've built}]
One more word for telling, explaining and describing at B1.
\end{grammarlens}

\subsection*{Your turn}
\begin{itemize}
\raggedright
  \item \emph{Say it:} \emph{créditer}
  \item \emph{Say it:} \emph{créditer}, once more
  \item \emph{Say it:} a sentence of your own with \emph{créditer}
  \item \emph{Recall:} say the French for cash, then the French for the withdrawal, then say \emph{créditer} again
\end{itemize}

\begin{tcolorbox}[breakable,colback=teal!4,colframe=teal!35!black,title={Before you move on}]
What does créditer mean? (\textquotedblleft{}To credit, to pay into\textquotedblright{}.) Say it once more.
\end{tcolorbox}
\section[régler]{régler — to settle (a bill), to pay; to adjust}

\label{lesson:FR-C451-regler}



\begin{quote}
Before the new one: say the French for the withdrawal, then the French for to credit.
\end{quote}

\subsection*{You'll want to know: régler}
\textbf{régler} — \textquotedblleft{}to settle (a bill), to pay; to adjust\textquotedblright{}.

\begin{quote}
Nous avons réglé la note de l'hôtel par carte.
\end{quote}

\begin{quote}
\emph{We paid the hotel bill by card.}
\end{quote}

\emph{Régler un problème} — to sort out a problem; \emph{régler la radio} — to tune the radio.

\begin{cousinweb}
\emph{régler} is built on \emph{règle} 'rule, ruler', which goes back to Latin \emph{rēgula} (EtymDB).
\end{cousinweb}

\begin{grammarlens}[title={What you've built}]
One more word for telling, explaining and describing at B1.
\end{grammarlens}

\subsection*{Your turn}
\begin{itemize}
\raggedright
  \item \emph{Say it:} \emph{régler}
  \item \emph{Say it:} \emph{régler}, once more
  \item \emph{Say it:} a sentence of your own with \emph{régler}
  \item \emph{Recall:} say the French for the withdrawal, then the French for to credit, then say \emph{régler} again
\end{itemize}

\begin{tcolorbox}[breakable,colback=teal!4,colframe=teal!35!black,title={Before you move on}]
What does régler mean? (\textquotedblleft{}To settle (a bill), to pay; to adjust\textquotedblright{}.) Say it once more.
\end{tcolorbox}
\section[le découvert]{le découvert — the overdraft}

\label{lesson:FR-C451-decouvert}



\begin{quote}
Before the new one: say the French for to credit, then the French for to settle.
\end{quote}

\subsection*{You'll want to know: le découvert}
\textbf{le découvert} — \textquotedblleft{}the overdraft\textquotedblright{}.

\begin{quote}
À la fin du mois, il est souvent à découvert.
\end{quote}

\begin{quote}
\emph{At the end of the month he is often overdrawn.}
\end{quote}

\emph{Être à découvert} — to be overdrawn.

\begin{cousinweb}
\emph{découvert} is the past participle of \emph{découvrir} 'to uncover' (EtymDB), from Late Latin \emph{discooperīre}: an account left uncovered by money.
\end{cousinweb}

\begin{grammarlens}[title={What you've built}]
One more word for telling, explaining and describing at B1.
\end{grammarlens}

\subsection*{Your turn}
\begin{itemize}
\raggedright
  \item \emph{Say it:} \emph{le découvert}
  \item \emph{Say it:} \emph{le découvert}, once more
  \item \emph{Say it:} a sentence of your own with \emph{le découvert}
  \item \emph{Recall:} say the French for to credit, then the French for to settle, then say \emph{le découvert} again
\end{itemize}

\begin{tcolorbox}[breakable,colback=teal!4,colframe=teal!35!black,title={Before you move on}]
What does le découvert mean? (\textquotedblleft{}The overdraft\textquotedblright{}.) Say it once more.
\end{tcolorbox}
\section[le relevé]{le relevé — the statement (from a bank); the reading}

\label{lesson:FR-C451-releve}



\begin{quote}
Before the new one: say the French for to settle, then the French for the overdraft.
\end{quote}

\subsection*{You'll want to know: le relevé}
\textbf{le relevé} — \textquotedblleft{}the statement (from a bank); the reading\textquotedblright{}.

\begin{quote}
Elle vérifie son relevé de compte chaque mois.
\end{quote}

\begin{quote}
\emph{She checks her bank statement every month.}
\end{quote}

\emph{Un relevé de compte}; \emph{le relevé du compteur} — the meter reading.

\begin{cousinweb}
\emph{relevé} is the past participle of \emph{relever} 'to lift; to note down' (EtymDB lists it), from Latin \emph{relevāre}, built on \emph{levāre} 'to raise'.
\end{cousinweb}

\begin{grammarlens}[title={What you've built}]
One more word for telling, explaining and describing at B1.
\end{grammarlens}

\subsection*{Your turn}
\begin{itemize}
\raggedright
  \item \emph{Say it:} \emph{le relevé}
  \item \emph{Say it:} \emph{le relevé}, once more
  \item \emph{Say it:} a sentence of your own with \emph{le relevé}
  \item \emph{Recall:} say the French for to settle, then the French for the overdraft, then say \emph{le relevé} again
\end{itemize}

\begin{tcolorbox}[breakable,colback=teal!4,colframe=teal!35!black,title={Before you move on}]
What does le relevé mean? (\textquotedblleft{}The statement (from a bank); the reading\textquotedblright{}.) Say it once more.
\end{tcolorbox}
